\section{Introduction}
\label{sec:intro}

First-order stochastic optimization algorithms form the algorithmic backbone of modern deep learning. Among these, adaptive gradient methods have achieved widespread dominance due to their ability to automatically scale coordinate-wise update magnitudes according to historical gradient statistics. The seminal work of \citet{kingma2014adam} introduced \textsc{Adam}, which synthesizes the advantages of \textsc{AdaGrad} \citep{duchi2011adaptive}—coordinate-wise scaling suitable for sparse gradients—and \textsc{RMSProp} \citep{tieleman2012rmsprop}—exponential moving averages capable of adapting to non-stationary objective landscapes. In conjunction with step-dependent bias correction for initialization transients, \textsc{Adam} rapidly became the default optimizer across computer vision, natural language processing, and reinforcement learning.

Despite its empirical ubiquity, the theoretical foundations of \textsc{Adam} contained a subtle yet fundamental flaw. In their breakthrough analysis, \citet{reddi2018convergence} demonstrated that the convergence proof of \textsc{Adam} in the standard Online Convex Optimization (OCO) framework \citep{hazan2016introduction} relies on an unfulfilled assumption: that the effective metric matrix $\Gamma_t = \frac{\sqrt{v_t}}{\alpha_t} - \frac{\sqrt{v_{t-1}}}{\alpha_{t-1}}$ remains positive semi-definite ($\Gamma_t \succeq 0$) for all iterations $t$. When objective functions feature infrequent but highly informative gradient bursts, the exponential moving average in \textsc{Adam} induces "short-term memory": the second-moment estimate $v_t$ decreases rapidly during non-burst iterations, causing the effective step size to explode in the wrong optimization direction. Consequently, \textsc{Adam} can suffer linear regret $\Omega(T)$ and diverge even on elementary one-dimensional convex objectives.

To resolve this defect, \citet{reddi2018convergence} proposed \textsc{AMSGrad}, which enforces monotonic non-decrease of the second-moment estimate via a running maximum: $\hat{v}_t = \max(\hat{v}_{t-1}, v_t)$. Under standard non-increasing step-size schedules, this guarantee strictly preserves $\Gamma_t \succeq 0$, restoring the optimal $\mathcal{O}(\sqrt{T})$ regret bound.

\paragraph{Contributions.} In this work, we present a complete, rigorous, and reproducible investigation of the \textsc{Adam} $\to$ \textsc{AMSGrad} transition:
\begin{enumerate}
    \item \textbf{From-Scratch Algorithmic Implementations:} We build a fully vectorized, framework-independent suite in pure NumPy comprising \textsc{SGD}, Momentum \textsc{SGD}, \textsc{AdaGrad}, \textsc{RMSProp}, \textsc{Adam}, and \textsc{AMSGrad}, verified by automated unit tests establishing exact $t=1$ debiasing and weak monotonicity.
    \item \textbf{Dual Counterexample Replication \& Diagnostic Instrumentation:} We replicate both the deterministic periodic construction (Theorem~3 of \citet{reddi2018convergence}) and the stochastic formulation (Section~3), introducing live tracking of $\Gamma_t$ to directly visualize the semi-definiteness violation.
    \item \textbf{Pre-Registered Phase-Boundary Extension:} We conduct an extensive parametric study across $(\beta_2, C)$ space with $N=100$ seeds per cell to evaluate the memory-horizon scaling law $\tau_2^* \propto T_{\text{burst}}$, reporting both the empirical collapse onto the dimensionless ratio $\rho = \frac{\tau_2}{T_{\text{burst}}}$ and the super-linear sensitivity of the divergence boundary.
\end{enumerate}
